%%=============================================================================
%% interviews
%%=============================================================================

\chapter{\IfLanguageName{dutch}{Interviews}{Interviews}}%
\label{ch:Interviews}


Een belangrijk deel van de bachelorproef bestaat uit het afnemen van 
interviews met professionals uit het werkveld. Hierdoor kan een compleet 
beeld gevormd worden van de competenties die nodig zijn voor een beginnende 
PL/1 en CICS developer op z/OS.

In dit hoofdstuk wordt toegelicht hoe de interviews zijn opgebouwd en welke 
vragen aan de professionals zijn gesteld. Het doel van deze interviews is dan 
ook niet enkel het verzamelen van verschillende meningen, maar het identificeren 
van terugkerende inzichten en het verkrijgen van een helder beeld van de 
exacte competenties die nodig zijn in het werkveld.


\section{Structuur en opbouw van de interviews}%
\label{sec:opbouw-van-de-interviews}

Voor de interviews werd er gekozen voor een semigestructureerde opbouw.
Aan de hand van deze methode is er een grotere flexibiliteit wat zorgt voor de voordelen van 
vaste vragen en voldoende ruimte om eigen accenten te leggen.
De flexibiliteit is ook belangrijk want het bepalen van competenties kan vaak niet gedaan worden door 
enkel het stellen van gesloten vragen maar heeft een bepaalde diepgang nodig die je in het gesprek kan bekomen.

De interviews zijn dan ook opgebouwd rond enkele kernvragen die rechtstreeks aansluiten bij de doelstellingen van de bachelorproef.
Deze doelstelling is het in kaart brengen van competenties die een startende developer nodig heeft om succesvol te kunnen starten in een mainframeomgeving.
Het stellen van deze vragen aan verschillende professionals in het werkveld kan ervoor zorgen dat er een terugkerend thema wordt gevonden.
Aan de hand hiervan kan men dan een actueel competentieframework opbouwen.

Hiervoor zijn volgende vragen gesteld aan de professionals binnen Euroclear:


\begin{itemize}
    \item \textbf{Vraag 1:} What slows down onboarding the most?
    
    Aan de hand van deze vraag kan achterhaald worden welke aspecten en factoren het integreren met 
    het team en opstarten van nieuwe werknemers vertragen. Ze geeft inzicht in de eerste problemen 
    waarmee starters geconfronteerd worden en toont aan waar de grootste frustraties liggen in de eerste
    weken in het werkveld.
    Hieruit kan dan bepaald worden welke competenties al aanwezig moeten zijn die het onboardingproces kunnen ondersteunen
    bij nieuwe developers.
    
    \item \textbf{Vraag 2:} What hard and soft skills do you wish juniors had before joining a mainframe team developing PL/1 and CICS?
   
    Met deze vraag werd er gepeild naar de technische en niet technische competenties die professionals graag aanwezig zien in
    junior developers. Uit de antwoorden kunnen we constateren welke voorkennis het meest waardevol is volgens professionals
    voor het starten binnen een PL/1 en CICS developmentteam.
    Deze skills kunnen dan ook dienen als een concrete basis voor het maken van het competentieframework en de proof of conceptsite.

    \item \textbf{Vraag 3:} What system-level knowledge does a starting PL/1 and CICS developer need?
    
    Aan de hand van deze vraag kan er een beeld geschetst worden van de nodige historische kennis.
    Daarnaast kan er ook getoetst worden naar welke andere tools en technologieën belangrijk zijn om een basis in te hebben
    om effectief te zijn in de rol van developer.
    
\end{itemize}

Met deze hoofdvragen als basis werden meerdere gesprekken gevoerd met professionals van Euroclear.
Elke professional legde eigen accenten afhankelijk van hun positie/specialisatie, maar er was wel een duidelijk patroon in de antwoorden.
Vooral het belang van globale systeemkennis en de historische achtergrond van de tools kwam steeds terug.
Daarnaast is een grondige kennis van de basisprincipes van programmeren niet te onderschatten in belang.


\section{Antwoorden van professionals}%
\label{sec:antwoorden-van-professionals}

Uit de interviews kan geconstateerd worden dat het onboardingproces vaak minder wordt vertraagd door het ontbreken van theoretische kennis,
maar door de complexiteit van de bestaande legacysystemen.
Beginnende developers kunnen vrij snel de basis programmeerconcepten oppikken, maar deze kennis toepassen in de bestaande systemen is veel moeilijker.
Hiervoor is kennis nodig van de historische groei van het systeem samen met de bedrijfsspecifieke manier van werken.

In het eerste interview met een CICS specialist Caio kwam dit aspect prominent naar voren.
Volgens hem is het essentieel dat een startende developer begrijpt waarom bepaalde technologieën vandaag nog bestaan en hoe deze zijn geëvolueerd.
Met deze kennis kunnen ze technologiekeuzes beter vatten en gaan ze sneller verbanden zien tussen oude en nieuwe systemen.

Volgens Caio kan men anders het complete beeld missen en denken dat iets verouderd is, ook al heeft het een heel belangrijke taak binnen het bedrijfsproces.
Maar deze kennis wordt ook grotendeels opgedaan in het bedrijf zelf, want overal gebruiken ze andere conventies en systemen.

Hierdoor haalde Caio een tweede punt aan dat volgens hem van belang is te beheersen als een starter.
Een algemene kennis van programmeerprincipes samen met kennis rond syntax.
Caio duidde erop dat het niet vereist is om expertkennis te hebben rond PL/1 en CICS,
maar dat men wel moet beschikken over een degelijke basis in programmeren.
Hier sprak hij van het begrijpen van variabelen, controlestructuren, foutafhandelingen en het lezen en begrijpen van bestaande code.

Die basis stelt een nieuwe starter in staat zich sneller aan te passen aan een specifieke omgeving.
Men kan dan ook sneller focussen op het leren van een diepere technische kennis, een die je enkel kan opbouwen door ermee te werken in een professionele omgeving.

Het tweede interview met een PL/1 expert Hervé bevestigde grotendeels de punten die Caio aankaartte.
Ook hij duidde aan dat uitgebreide systeemkennis niet vooraf wordt verwacht dat die aanwezig is.
Dit omdat het volgens Hervé het aanleren van dit meer iets is dat door ervaring in het bedrijf komt dan door zelfstandige studie.
Wat hij wel belangrijk of zelfs als noodzakelijk zag was een sterke basis in programmeerprincipes, en het vertrouwd zijn met de syntax van PL/I.
Volgens Hervé moet een starter wel in staat zijn om code te kunnen analyseren en bepaalde logische structuren te herkennen.

Naast het analyseren is het debuggen van code volgens Hervé misschien wel even belangrijk.
Wie in staat is om foutmeldingen correct te analyseren en stap per stap jobuitvoer kan begrijpen, die starter zal in staat zijn om zelfstandig te leren uit eigen fouten.
Debuggen is ook iets dat even belangrijk blijft, een ervaren coder gebruikt deze skill nog steeds dagelijks.


\section{Conclusie van de interviews}%
\label{sec:Conclusie-van-de-interviews}

Volgens de professionals is het dus niet realistisch om van een junior te verwachten dat die van dag één alle systeemkennis van het bedrijf kent.
Veel van die kennis is gekoppeld aan hun specifieke bedrijfsomgeving en is iets dat aangeleerd wordt naarmate men meer reële toepassingen uitwerkt.

Maar een concrete basis in programmeerprincipes en kennis van de tools is een must.
Het bezitten van deze basis stelt de starter in staat om zelfstandig bedrijfsspecifieke systeemkennis te leren.